\section{Related Work}
\label{sec:related}

\subsection{Agentic and Adaptive RAG}
Self-RAG~\cite{selfrag2024} introduced reflective retrieval tokens.
CRAG~\cite{crag2024} added a corrective lightweight retrieval evaluator.
Adaptive-RAG~\cite{adaptiverag2024} routes among no-retrieval,
single-step, and multi-step retrieval based on a classifier-predicted
query complexity. Search-R1~\cite{searchr1_2025} trains the retrieving
agent to interleave search and reasoning via reinforcement learning.
RAG-Critic~\cite{ragcritic2025} adds an agentic critic loop with
iterative refinement. None of these systems condition routing on
typed-graph distance, and none expose a typed graph-lookup tool;
retrieval is treated as flat text top-$k$.

\subsection{GraphRAG Family}
Microsoft GraphRAG~\cite{graphrag2024} builds entity--relation graphs
and answers via community summaries; LightRAG~\cite{lightrag2024} and
HippoRAG~2~\cite{hipporag2_2025} reduce graph-construction cost.
Han et~al.~\cite{ragvsgraphrag2025} compare RAG and GraphRAG empirically;
GraphRAG-Bench~\cite{graphragbench2026} provides a benchmark for
\textit{when} graph structure helps. Graph-R1~\cite{graphr1_2025} trains
a graph-aware retrieval policy with RL. These systems exploit graph
structure but are non-agentic: they do not loop a critic over retrieved
evidence, and they do not condition pipeline choice on hop distance.

\subsection{RAG for Requirements Traceability}
LiSSA~\cite{lissa2025} applies single-shot RAG to source--requirement
traceability. TraceLLM~\cite{tracellm2026} and TVR~\cite{niu_tvr2025}
extend LLM-based traceability with verification feedback. Graph-RAG
for compliance~\cite{graphragcompliance2024} uses domain knowledge
graphs without an agentic loop. To our knowledge, this is the first
work to combine an agentic critic loop, a typed-edge graph tool, and
hop-stratified evaluation in the requirements-traceability setting.
